\section{Síntesis y ampliación}

\subsection{Un hilo conductor que atraviesa toda la clase: controlar las atajos que el modelo puede ver}

Dos trabajos parecen tratar temas distintos, uno sobre objetos y otro sobre regularización gaussiana, pero en realidad ambos rediseñan el espacio de soluciones factibles. Causal-JEPA elimina el historial del propio objeto objetivo, haciendo necesaria la información de relaciones fuera de las atajos de interacción; LeWorldModel restringe la distribución de embeddings, impidiendo que una representación constante sea la solución óptima legítima para la pérdida de predicción. El principio común del aprendizaje autosupervisado de alta calidad es hacer difícil satisfacer el objetivo con soluciones ``baratas pero incorrectas''.

\begin{importantbox}{Las ocho preguntas más transferibles de esta clase}
\begin{enumerate}
  \item ¿Qué variables relevantes para la decisión conserva la representación del estado?
  \item ¿Cuál es el atajo más fácil de explotar?
  \item ¿Qué soluciones triviales excluye la restricción anti-collapse?
  \item ¿Por qué interfaz ingresa la acción al predictor?
  \item ¿La información visible durante el entrenamiento coincide con la del momento de inferencia?
  \item ¿El aumento de resultados proviene de la representación, la arquitectura, la función de pérdida o el controlador?
  \item ¿Qué no pueden probar por sí solos las sondeadoras (probe), la atención y los despliegues bonitos (rollout)?
  \item ¿Cómo fallan con horizontes largos, acciones fuera de distribución y entornos reales?
\end{enumerate}
\end{importantbox}

\subsection{El sesgo estructural y el sesgo de optimización deben trabajar en conjunto}

Solo con ocultamiento de objetos (object masking) y sin una representación estable, el modelo podría aprender supuestas interacciones sobre latentes degenerados; solo con SIGReg y sin estructura de tarea, el modelo podría mantener una varianza rica pero seguir dependiendo de correlaciones de fondo. Los modelos del mundo de la próxima generación deben diseñar simultáneamente unidades de representación, intervenciones de observabilidad, distribuciones anti-collapse, interfaces de acción y planificadores. Todas ellas deben ser ablatadas individualmente dentro del mismo protocolo de evidencia.

\begin{knowledgebox}{La relación complementaria entre Causal-JEPA y LeWorldModel}
Causal-JEPA responde principalmente ``¿qué variables debería depender la predicción?''; LeWorldModel responde principalmente ``¿cómo mantiene el encoder información?''. El primero actualmente depende de un encoder centrado en objetos congelado, mientras que el segundo utiliza latentes compactos globales. Una extensión natural sería aprender de manera end-to-end slots de objetos estables, manteniendo al mismo tiempo las máscaras que fuerzan la interacción y regularizadores anti-collapse escalables.
\end{knowledgebox}

\subsection{Tres límites de evidencia}

Primero, los contramodelos (counterfactual), la planificación y las ganancias por sorpresa en benchmarks controlados no equivalen a una comprensión causal del mundo abierto. Segundo, que un latente sea decodificable y que la atención esté concentrada no significa que el predictor utilice realmente el mecanismo correcto. Tercero, los títulos de los artículos y las formulaciones teóricas en el momento de la clase están sujetos a revisión por pares y análisis posteriores; las apuntes de investigación deben guardar instantáneas históricas y también marcar los límites más prudentes establecidos posteriormente.

\begin{warningbox}{No escribas modelos del mundo como una capa intermedia omnipotente}
Los modelos del mundo comprimen las observaciones; la compresión implica pérdida de información. Pueden generar despliegues (rollout), lo que acumula errores; pueden ser llamados por un planificador, quien podría explotar el error del modelo. Son una interfaz potente que conecta predicción, control y verificación, no una capa mágica que elimine automáticamente la alucinación, la seguridad y el cambio de distribución.
\end{warningbox}

\subsection{Preguntas abiertas}

Las direcciones que vale la pena seguir incluyen: cómo descubrir y rastrear objetos de manera estable en videos reales; cómo incorporar múltiples escalas de tiempo y abstracciones de eventos en los latentes; cómo marcar con incertidumbre el futuro que el modelo no conoce; cómo hacer coexistir objetivos de lenguaje, recompensas, restricciones y objetivos de imagen; cómo evitar que el planificador explote el error del modelo; y cómo distinguir correlaciones, suficiencia predictiva y mecanismos causales identificables mediante intervención con robots reales.

\subsection{Lecturas adicionales}

Se sugiere leer Causal-JEPA v1 y LeWorldModel v1 según el momento de la clase, luego contrastarlos con las revisiones posteriores de Causal-JEPA para observar cómo el lenguaje causal se vuelve más cauteloso. El material contextual puede leerse secuencialmente: Atención a Slots (Slot Attention), V-JEPA, V-JEPA 2 y DINO-WM. A nivel de implementación, pueden consultarse los códigos oficiales de ambos trabajos y la librería de experimentos unificada `stable-worldmodel`. Al leer, continúe utilizando las ocho preguntas de esta clase en lugar de comparar solo cifras del leaderboard.

\end{document}
